\documentclass[letterpaper,11pt]{article}
\usepackage[empty]{fullpage}
\usepackage[hidelinks]{hyperref}
\usepackage[english]{babel}
\addtolength{\oddsidemargin}{-0.6in}
\addtolength{\evensidemargin}{-0.5in}
\addtolength{\textwidth}{1.19in}
\addtolength{\topmargin}{-.9in}
\addtolength{\textheight}{1.7in}
\urlstyle{same}
\raggedbottom
\raggedright
\setlength{\parindent}{0pt}
\setlength{\parskip}{10pt}
\pagestyle{empty}

\begin{document}

\begin{center}
    {\LARGE \scshape Jendrick Rian Galut} \\ \vspace{2pt}
    \small Surrey, BC ~$|$~
    672-272-3193 ~$|$~
    \href{mailto:rcgalut1@gmail.com}{\underline{rcgalut1@gmail.com}} ~$|$~
    \href{https://www.linkedin.com/in/jendrick-rian-galut-374310307/}{\underline{linkedin.com/in/jendrick-galut}} ~$|$~
    \href{https://github.com/Jeam-cloud}{\underline{github.com/jeam-cloud}}
\end{center}

\vspace{8pt}

August 12, 2026

Hiring Manager \\
KPMG in Canada \\
Technology Risk Services

Dear Hiring Manager,

I am writing to apply for the Summer 2027 Technology Risk Services Summer Intern position with KPMG in British Columbia. I am a third-year Computing Science student (Business minor) at Trinity Western University, and I am drawn to this role because it sits exactly where I already spend my time: making sure a system's access controls, audit trails, and safeguards actually hold up under scrutiny.

That interest is not abstract. I am currently building ARTHUR, a local-first AI desktop assistant, and the hardest part of that project has been risk control, not features: I designed a task-mode privilege separation model and an approval broker that gates any risky action behind explicit confirmation, sandboxed code execution in Docker with per-tool network policy, and added a static-analysis scan surfaced directly in the approval dialog. When I identified a prompt-injection exposure in how the assistant handled untrusted web content and file attachments, I mitigated it by spotlighting that content between explicit markers before it reaches the model, and verified the fix with a dedicated pytest suite rather than taking it on faith. That is, in miniature, the discipline I understand Technology Risk Services to practice at client scale: identify where a system is exposed, design a control, and prove the control works.

I bring a similar instinct for controls and accountability to my work outside the classroom. As a Building Manager, I administer access-control and identity-credentialing systems for over 100 residents, track every incident and bylaw issue through a structured case-management system, and maintain digital audit trails across the property. As a Teaching Assistant for two computing science courses, I have graded against consistent rubrics and delivered structured, evidence-based feedback to more than 25 students each term, which has sharpened how clearly I can explain a technical finding to someone who was not in the room when it happened.

I would welcome the chance to bring that combination of technical grounding and attention to controls to KPMG's Technology Risk Services team this summer. Thank you for your time and consideration, and I look forward to the opportunity to discuss my application further.

Sincerely,

Jendrick Rian Galut

\end{document}
